\documentclass[12pt,reqno]{amsart}
\usepackage{../handout, ../customcommands}
\usepackage{caption}

\hdr{MAT 240}{\S2.1 The definition of conditional probability, part 1}

\begin{document}

\textbf{Suggested problems:} 1-5 

\bigskip
\hrule

\bigskip

\prob (Example 2.1.1) Consider a state lottery game in which six numbers are drawn without replacement from a bin containing the numbers $1$-$30$. Each player tries to match the set of six numbers that will be drawn without regard to the order in which the numbers are drawn. Suppose that you hold a ticket in such a lottery with the numbers $1$, $14$, $15$, $20$, $23$, and $27$. You turn on your television to watch the drawing but all you see is one number, $15$, being drawn when the power suddenly goes off in your house. You don’t even know whether $15$ was the first, last, or some in-between draw. However, now that you know that $15$ appears in the winning draw, the probability that your ticket is a winner must be higher than it was before you saw the draw. Calculate the revised probability of your ticket is the winning one.

\vfill
\prob (Example 2.1.3) Suppose that two dice were rolled and it was observed that the sum $T$ of the two numbers was odd. Determine the probability that $T$ was less than $8$.

\vfill
\newpage
\prob
Suppose that two balls are to be selected at random, without replacement, from a box containing $r$ red balls and $b$ blue balls. Determine the probability $p$ that the first ball will be red and the second ball will be blue.

\vfill
\prob Suppose that four balls are selected one at a time, without replacement, from a box containing $r$ red balls and $b$ blue balls ($r\geq 2$, $r\geq 2$). Determine the probability of obtaining the sequence of outcomes red, blue, red, blue.

\vfill
\end{document}